\subsection{RQ2: Is it the horizon or the algorithm?}
\label{sec:rq2}

PPO differs from the bandit in two ways at once: it optimizes a discounted
return with $\gamma=0.995$, and it is an on-policy policy-gradient method
rather than a value regressor on replayed data. We vary the discount within
each family on the three preregistered roots, holding everything else fixed
(Table~\ref{tab:horizon}, Fig.~\ref{fig:curves}).

\emph{Within PPO}, removing the horizon closes the gap. PPO with $\gamma=0$
reaches \HPPOZero{} against the bandit's \HQZero{} on the same roots (mean
difference \HPPOZeroVs, above the bandit on \HPPOZeroRootsBetter{} of
\HPPOZeroRoots{} roots), while PPO with $\gamma=0.995$ reaches \HPPODefault{}
(above the bandit on \HPPODefaultRootsBetter{} of \HPPODefaultRoots{} roots). The oscillation disappears with the horizon: reversals
fall from \HPPODefaultRev{} to \HPPOZeroRev{} per episode, and the validation
curves become monotone (Fig.~\ref{fig:curves}). This points to the long horizon,
rather than the policy-gradient algorithm, as what limits PPO here. It does not
exclude an interaction between PPO and long horizons---for example the scale of
undiscounted-sum value targets---which the reward-normalization and GAE
configurations of \S\ref{sec:rq3} probe.

\emph{Within the value family}, adding a horizon does not help. The masked
Q-learner is the bandit with one change---a bootstrapped, discounted target
(Eq.~\ref{eq:qtarget}). With $\gamma=0.5$ it is close to the
bandit (mean difference \HQPointFiveVs, better on \HQPointFiveRootsBetter{} of
\HQPointFiveRoots{} roots); with $\gamma=0.9$ it is worse (\HQPointNine{} vs
\HQZero; mean difference \HQPointNineVs, better than the bandit on
\HQPointNineRootsBetter{} of \HQPointNineRoots{} roots), mainly through lower
network utility rather than higher churn (\HQPointNineRev{} reversals per
episode); with $\gamma=0.99$ it falls to \HQPointNineNine{} (mean difference
\HQPointNineNineVs, better on \HQPointNineNineRootsBetter{} of
\HQPointNineNineRoots{} roots). In both families the return decreases as the
discount grows beyond $0.5$. This matches the sequentiality diagnostic: the part of the future
that a learner could exploit from its observation is small, so a bootstrapped
target adds estimation error without adding much signal---the effect of a
shorter planning horizon under limited data~\citep{jiang2015horizon,amit2020discount}.

\emph{Shaping.} The training reward contains a potential-based shaping term
that is policy-invariant only for $\gamma=0.995$, so it could distort the
comparison across discounts. Removing it on two roots (42, 314159; evaluated
on the primary reward) leaves the ordering intact.
\ifcsname NoShapePPOVsBandit\endcsname Relative to the shaping-free bandit,
shaping-free PPO with $\gamma=0$ differs by \NoShapePPOZeroVsBanditPerRoot{},
the Q-learner with $\gamma=0.9$ by \NoShapeQVsBanditPerRoot{} and PPO with
$\gamma=0.995$ by \NoShapePPOVsBanditPerRoot{} (roots 42 / 314159). The
differences that survive learner noise (\S\ref{sec:rq3}) are the large ones,
and they point the same way with and without shaping.\fi


\begin{table}[t]
\centering\small
\caption{Horizon sweep (test return, selected checkpoints) on up to three
training roots (42, 314159, 271828; column Roots). The value learner with
$\gamma=0$ is the bandit. ``vs.\ bandit'': mean paired difference to the
bandit on the same roots, and per root in the order listed.}
\label{tab:horizon}
\setlength{\tabcolsep}{3.5pt}
\begin{adjustbox}{max width=\columnwidth}
\begin{tabular}{@{}lrrrrrr@{}}
\toprule
Learner & $\gamma$ & Roots & Return & vs.\ bandit & Rer./h & Rev. \\
\midrule
\tabinput{tables/horizon_sweep}
\bottomrule
\end{tabular}
\end{adjustbox}
\end{table}

\begin{figure}[t]
\centering
\includegraphics[width=\columnwidth]{fig_learning_curves}
\caption{Mean validation return (seeds 2001--2005) during training. Thin lines:
individual roots. PPO with $\gamma=0.995$ is unstable across roots; with
$\gamma=0$ it learns as fast as the bandit.}
\label{fig:curves}
\end{figure}
